\section{Setting and Technical Tools}
In Online Packet Scheduling with Deadlines (OPSD), an instance is a sequence of packets $\mathcal{P} = \{(r_{p},d_{p}, w_p)\}_{p}$, where $r_p$ and $d_p \ge r_p$ are integers representing the arrival time and deadline of the packet $p$, respectively. $w_p\ge 0$ is a real number that represents the \textit{weight} of packet $p$. For the rest of this paper, we will slightly abuse notation by calling $w_p$ simply $p$. The time is discrete and divided into slots. A \textit{schedule} assigns a subset of packets $\mathcal{S} \subset \mathcal{P}$ to time slots such that (i) any packet $p \in \mathcal{S}$ is assigned to a slot in $\{r_p, \ldots, d_p\}$ and (ii) each slot has at most one packet assigned. The objective is to design a schedule that maximizes the total weight of the packets in $\mathcal{S}$. In the \textit{$s$-bounded} variant of the problem, the set of instances is constrained to those where $d_p \le r_p+s-1$, \emph{i.e.}, every packet can be assigned in either the arrival slot or one of the next $s-1$ ones (for a total of $s$ possibilities). We call $T \coloneqq \max_{p\in \mathcal{P}} d_p$. An online algorithm \alg observes the packets arriving over time and, at every round $t \in [T]$, has to decide which packet to schedule without observing the future. The packets not scheduled at time $t$ are stored in a set $\mathcal{B}_t$ called \textit{buffer}, until their deadline.
\paragraph{Competitive Ratio } We call $\Galg$ the total gain of \alg, \emph{i.e.} the total sum of the packets scheduled up to $T$.
The standard performance metric in the literature is the \textit{competitive ratio}, \emph{i.e.}, the ratio between the performance of a benchmark decision-maker and $\Galg$ \citep{borodin2005online}. The goal is to minimize the competitive ratio. The adversary observes the algorithm and decides the input sequence. If the algorithm is deterministic, the full sequence can be chosen beforehand, as the behavior of the algorithm is completely predictable. In the case of randomized algorithms, the adversary cannot observe its random bits during the execution, but can decide the next entry in the sequence based on the past decisions of the algorithm. The benchmark is the best possible decision-maker that observes the whole sequence, called \opt. Thus, the competitive ratio $C$ is defined by the inequality $\Gopt \le C\cdot\Galg$. 

% As customary in the literature, we distinguish between deterministic and randomized algorithms. Randomized algorithms can attain strictly better competitive ratios than deterministic algorithms in every setting. However, the theory behind randomized algorithms is less explored than its deterministic counterpart. In particular, in the general case, there still is no algorithm matching the competitive ratio lower bound of $\frac{5}{4}$. The literature on $2$-bounded instances is the largest so far. This special case is practically relevant, and is not worst-case easier than the general scenario (the lower bound of $\Phi$ from \cite{hajek2001competitiveness} is proven on a $2$-bounded instance). In this paper, we place a particular focus on $2$-bounded instances, and show that learning can be done with tight guarantees in this special scenario. 

\subsection{Learning in Online Packet Scheduling with Deadlines} 
\label{sec:prob_form}
We consider the \textit{Stochastic Bounded-Delay $K$-Packet Scheduling} problem, a variant of the problem where there exist $K$ distinct classes of packets, and all packets belonging to the same class have their weights sampled from the same probability distribution with mean $\mu_k$ for $k\in[K]$. We then define an instance as $\mathcal{P}= \{(r_p, d_p, X_p, c_p)\}_{p}$, where $c_p \in [K]$ indicates the class of the packet $p$ and $X_p$ is the weight of packet $p$, which is sampled from the distribution $\mathcal{D}_{c_p}$ with mean $w_{p} \coloneqq \mu_{c_p}$. Without loss of generality, we order the expectations of the $K$ types in a descending order, \emph{i.e.}, $\mu_1 > \mu_2 > \ldots > \mu_K$. Moreover, we define $\Delta_{i,j} \coloneqq \mu_i - \mu_j$ for $j \ge i$ as the \textit{sub-optimality gap} of $j$ w.r.t. $i$. We will slightly abuse the notation and indicate $w_p$ simply as $p$ for the rest of this paper. The scheduler does not observe the distributions nor their means, and observes $X_{p} \in [0,1]$ at time $t \in [T]$ only if the packet $p$ is assigned to slot $t$. The adversary is allowed to decide the type of the packet injected, but not the realization of the weight (and does not observe it beforehand). This variant can be trivially extended to the $s$-bounded setting. As it will become handy in what follows, we introduce the set $V_t \subseteq [K]$ as the set of packet types such that at least one packet of the type is in the buffer and has a deadline $t$. We call $B_t \subseteq [K]$ the set of packet types that are not in $V_t$ and such that at least one packet of the type is in the buffer and has a deadline \textit{greater than} $t$. Note that no reasonable \alg schedules a packet of type $i$ with deadline greater than $t$ if $i \in V_t$.

\paragraph{$\alpha$-Regret} First, consider a deterministic algorithm \alg. Let $\Galg = \sum_{i\in \mathcal{S}} X_i$ be the total gain of \alg. We assume that no adversary can observe the actual realizations of the weights, and consider an \opt that does the best possible scheduling in expectation, \emph{i.e.} uses the weight's expectation as the actual value. We call $\mathbb{E}[\Gopt]$ the expected total gain of \opt. Thus, a natural performance metric is the \textit{$\alpha$-regret}, defined as $R_{\alpha,T} = \mathbb{E}[\Gopt] - \alpha\mathbb{E}[\Galg]$, where the expectation is taken w.r.t. the weights' randomization (that is, in this case, the only possible source of randomness). If \alg is randomized, the expected total gain also accounts for the algorithm's randomization. The definition of \opt is the same: \opt can always be assumed deterministic as the provided sequence is fixed. Finally, we say that an algorithm achieves \textit{$\alpha$-no-regret} if $R_{\alpha,T}^{\text{\alg} }= o(T)$, \emph{i.e.} its $\alpha$-regret is sublinear in $T$.

\paragraph{Deterministic vs Randomized Algorithms in $K$-OPSD}
In $K$-OPSD we need to set a formal definition of a deterministic algorithm. Indeed, due to the stochastic nature of the weights, the algorithm's decision is intrinsically stochastic. Thus, as customary in the bandits literature, we say that an algorithm is deterministic if its decision at time $t$ is \textit{predictable} w.r.t. the filtration up to time $t-1$. In other words, the adversary can predict which packet is processed by \alg at time $t$ by just observing the history until $t-1$. For instance, the well-known \texttt{UCB1} algorithm \citep{auer2002finite} is deterministic in this sense. One can interpret the distinction between deterministic and randomized algorithms as a distinction between the strength of the adversary: while competitive ratios are lower in randomized algorithms, it is also true that the adversary is less informed. An adversary observing the actual random bits of the algorithm would result in every algorithm being, \textit{de facto}, deterministic.

\paragraph{Mean Estimation and Confidence Bounds}
Define $\widehat{\mu}_{i,t} = \frac{1}{N_i(t)} \sum_{s=1}^t r_i \mathds{1}_{p_s = i}$ to be the empirical mean of the observed weights from packets of type $i$ observed up to time $t$. Let $\mathrm{LCB}_{i,t}(\delta)$ and $\mathrm{UCB}_{i,t}(\delta)$ be lower and upper confidence bounds on the unknown mean weight from a packet of type $i$, with confidence level $1-\delta$.  Let $\beta_{i,t}(\delta) \coloneqq \frac{1}{2}(\mathrm{UCB}_{i,t}(\delta) - \mathrm{LCB}_{i,t}(\delta))$ be (half of) the width of the confidence bound. We use the same procedure introduced by~\cite{levy2024online} to construct our confidence bounds. Specifically, the confidence width is chosen to be $\beta_{i,t}(\delta) = \sqrt{\frac{\log(K T^2/\delta)}{2 N_i(t)}}$, where $K$ is the number of bounds that must hold simultaneously (in our case, one per packet type). Then, $\mathrm{UCB}_{i,t}(\delta) \triangleq \min\{ \mathrm{UCB}_{i,t - 1}(\delta), \widehat{\mu}_{i,t} +\beta_{i,t}(\delta) \}$, and similarly, $\mathrm{LCB}_{i,t}(\delta) \triangleq \max\{ \mathrm{LCB}_{i,t - 1}(\delta), \widehat{\mu}_{i,t} - \beta_{i,t}(\delta) \}$. Lemma C.1 in~\cite{levy2024online} ensures that the true weight lies inside the confidence bound with probability at least $1 - \frac{1}{KT}$ for a proper choice of $\delta$, and that the UCBs and LCBs are positive and monotonic. For clarity, we will omit the dependence on $\delta$ (except when necessary) and assume that $\delta \le \frac{1}{T}$. $\delta$ is an algorithm parameter (it is necessary to compute UCBs and LCBs), and additional details on the choice of $\delta$ for every algorithm can be found in the proofs.


\subsection{Connection with Multi-Armed Bandits}
\label{sec:sleeping}
In the Sleeping bandit problem~\citep{kleinberg2010regret}, a learner is sequentially faced with an \emph{adversarially} chosen set of actions $\mathcal{A}_t \subset [K]$ for $T$ rounds. Each action is associated with a probability distribution that the learner does not observe, and a sample (or \textit{reward}) is received after the corresponding action is chosen. The cumulative reward of the learner is given by the sum of all rewards obtained up to $T$. This problem, which is a generalization of the stochastic $K$-armed bandit problem, has been explored in previous literature under the lens of \textit{regret minimization}. In regret minimization, the learner's goal is to minimize the performance gap w.r.t. the best possible learner, \emph{i.e.} $R_T \coloneqq \mathbb{E}[\Gopt]-\mathbb{E}[\Galg]$ \cite{lattimore2020bandit}. Note that regret minimization is a special case of $\alpha$-regret minimization ($\alpha = 1$). As the sleeping bandit problem includes the stochastic $K$-armed bandit problem, the lower bound on the regret that any algorithm can suffer is greater or equal, and it is in the order of $\Omega ( \sqrt{KT})$. Interestingly, the $K$-OPSD problem is a generalization of the \textit{$K$-armed} sleeping bandit learning problem. 
\begin{restatable}[Sleeping Bandits are a special case of $K$-OPSD]{lemma}{reductionMABS}
Every $1$-bounded instance of the $K$-OPSD problem can be reduced to a $K$-armed Sleeping Bandit problem, and vice versa.
\end{restatable}
The \texttt{AUER} algorithm \citep{kleinberg2010regret}, a gold standard for regret minimization in sleeping bandits, employs the well-known \textit{optimism in the face of uncertainty} principle \citep{auer2002finite} to achieve no-regret. \texttt{AUER} suffers an expected regret upper bounded by $\widetilde{\mathcal{O}}\left(\sqrt{KT}\right)$. 